\section{Puzzles}
\begin{quotation}\textit{
There are some computing challenges that exist almost solely because
their solutions are not immediatly obvious, or as puzzles, jokes,
tests to set job candidates or purely because a challengs can be fun.
But some of these can contain lessons that might matter in the real
world. The first one considered here is a classical challenge that dates
back to the very early years of computing, but after explaining what is
called for and how to achieve the effect there is going to be a discussion
of a few ways in which this has importance way beyond the frivolous joke it
starts off as.
}\end{quotation}

\subsection{Quine}
The absurd seeming challenge is to create a program wuch that when it is run
it displays exactly itself in the textual source-code form it started in.
It should instantly be obvious that it can not just consist of one or
more print statements to generate this, because the full code will of
necessity be bulkier than the strings that are to be printed. The exact
programming language to be used here will make a bit of a difference
to exactly what is to be done, so here a form of pseudo-code that is
not exactly any existing language will be used -- but conversion to
a version that could be demonstrated should be easy.

A self-reproducing program like this is referred to as a ``quine'' in
honour of Willard Van Orman Quine, a philosopher who had much to say
about the roots of mathematics, including issues relevant to such
a paradoxical piece of code, but who did not himself promote the
particular challenge.

The insight needed to solve the challenge comes from the fact that
to get enought output to be a complete copy of the source text, a
program must display some of its data (at least) twice. So one
can give a high level skeleton of a solution as
\begin{verbatim}
start here:
   let d = "some useful text that we will design in a moment";
   first_display(d);
   second_display(d);
   third_display(d);
   stop
where
     first_display(s) = ...
and  second_display(s) = ...
and  third_display(s) = ...
\end{verbatim}
Where it is now necessary to fill in both the text that is set up in
the variable \verb+d+ and the details of the two or three functions
that display it.

Well it is easy to get started because \verb+first_display()+ must
begin by printing out something along the lines of
\verb+start here: let d=+ and so if we make the string d start with
those characters it can just print out the first so many chartacters
from its argument. What comes next has to be the material that
forms the totality of \verb+d+. So \verb+second_display+ just pushes
out the whole of what is in \verb+d+. Well it is liable to have to
take special action to render quote marks and line breaks and where
they are to be needed for any reason the string may have to use
some encoding for them -- bit it can be rather simple encoding.
Finally \verb+third_display()+ prints out the parts of the entire program
from the end of the definition of \verb+d+ onwards, in very much the
same way that \verb+first_display()+ showed the start.
The effect is that the source code has effectively two versions of everything.
one as executable code and the second as part of the data that it decodes
to know what to print. Once you have seen the trick it is not really
that difficult after all!

However the idea is not merely an abstract puzzle. I can form the
basis of a way to create a nice form of malicious software, and it
is central to allowing proof of some results in the theory of computation.

To understand the malicious use consider a compiler for some programming
language. You can feed it source code and it generates executable
binaries. Quite often compilers will be written in the language that they
process. making the compiler self-reliant like this is sometimes referred to
as ``eating your od dog-food'' or just ``dogfooding'', and it can be
good policy because it can represent a fair-sized test to check that
the compiler works. Now imagine that the initially shipped binary includes
an extra section that inspects the code it is asked to compile and takes
special action if it matches certain signatures. One part of the special
action will be to feed the text of this special extra code into the
source that it works on even if it had not been present in the source
files that had been offered. This is very much like using a quine. It
means that even when presented with clean sources of a respectable version
of itself it will generate an executable containing its cunningly hidden
trap. As well as using this to hide the fact that it is a ``trick'' version
of the code, it can also include (in the hidden part) tests that detect
features it wishes to attack. This could be as simple as making it
inspect copyright notices and then arranging that executables it builds
from code owned by some person or company includes extra material (again
somewhat in the spirit of a quine) that arranges to post the source
code (which might well be supposed to be confidential) back home.
This shows that installing the binary of a compiler from anywhere outside
your direct control could carry risks. And that to mitigate those you
may need to rebuild it from its source (which is assumed human readable and
visibly untainted) using some independent set of build tools. Most people
are not neurotic enough to feel the need to go to this level of trouble,
but if company or national security is at stake it may need to be considered!

The second use of code that in effect contains a copy of itself is at the
heart of one of the more important theoretical results about what computers
can and can not do. This result asserts that it is impossible to have a
program that guarantees to judge whether some code you show it will
ever terminate or whether it will run for ever. This is the Halting Problem.

A pedantic presentation of the Halting Problem has to pick some particular
concrete model of what computation is and what notation the program that is
to be assessed must be presented in. There needs to be consideration as to
whether some programming languages or models of computation sidestep the
problem either when used to code up the termiantion-tester or when the
code to be tested is expressed using them. The view taken here is that these
are all important technical issues that must be addressed in a formal proof,
but that it can help to come to terms with a much more informal explanation
first, especially since that can be kept quite compact and simple.

So to show that it is not possible to have a program such that if you
pass it some source code and some data, it will tell you if the
computation descibed by that source code as applied to the data will
ever finish. There are of course two caveats here. The first is that
for many examples of the source code (i.\e.\ all those corresponding to
``sensible'' programs) it will be possible to analyse what can go on and
judge termination. The second is that it would of course be possible to
execute or simulate the behaviour as described by the source code and data
and if the calculation described terminates the simulation will complete
and it is possible to report ``yes this one does eventually stop''. The
problem at the core of the Halting Problem is that it turns out there is
no way of being certain that you can distinguish bewteen ``it has not
finished yet'' and ``it is never going to''.

The following example program is {\em not} one that illustrates undecidability
but is one that suggests that judging termination may be a seriously tough
challenge: In about 1637 Fermat asserted that there could be no nontrivial
whole number solutions to the equation $x^N + y^N = z^N$ for $N >= 3$. For
aound 350 years nobody could show that this was the case, so there might
have been integers $x$, $y$ and $z$ that satisfied the equation, although over
time it was established that if they existed at all they would have to be
gigantic. It was only in 1994-1995 that Andrew Wiles presented a proof
that Fermat's assertion had been correct. So until then a program
of the form
\begin{verbatim}
    for K = 2, 3, ...
       for all possible x, y, z, N each in the range 1:K with N>=3
          if x^N + y^N = z^N stop
\end{verbatim}
was one where nobody could be certain as to whether the stop statement would
ever be activated or whether the program would run for ever. Of course this
is the sort of program that is aimed at theoreticians -- it supposes that
there is no limit to the magnitude of the numbers that arise and it is
happy to look for a counter-example to Fermat's claim using a really crude
exhaustive search. Howver it is a short and fairly transparent program
where it took the world's mathematicians 350 years before they could be
certain that it would never terminate. Programs that can be hard to analyse
do not have to be unduly long!

Now we will show how to prove the Halting Problem can not be solved in
general. The key idea is to start by pretend that it could be solved and
we have the program that solves it. It is then possible to show that the
very existence of this program leads to an absurdity -- in some manner
via a quine. Here one wants a sight variation on the idea of applying the
testing program to its own source code with another copy of its own source
code as the data that is to be presented to the version under test. The way
that this is self-referential cab be achieved in very much the way that the
quine program included enough of a copy of itself in its data string. So it
is mildly tricky but managable.

Well what is actually needed is a slight variant on code that tests a
program and prints either ``terminates`` or ``runs for ever''. If
this adjusted version judges that the program that it is testing it
will merely go into a simple unending loop, and so it itself will not
terminated. If on the other hand it concludes that the program under test
is going to run for ever it just halts. These changes are really rather minor
and are clearly trivial to implement.

Based on what we now feed the program a copy of itself trying to
analyse itself. If we really had a program that has solved the Halting Problem
then the outer version of the code would be able to confirm that the
version it was testing was going to terminate. But then it goes into
that infinite loop and does not complete, so it is unable to report
the success. On the other hand if the version under test would loop
for ever the outer one would spot that and halt very happily. This situation
represents an intolerable situation where what we have will halt if it
does not half, or loop if it does helt. Something along the line has
gone badly wrong! The only supposition we made was that there was a program
that solved the Halting Problem, so we have to abandon that hope and conclude
that there can never be one. It is probably obvious that a fundamental
limit on what computers can ever do can have a lot of impact.


As a third way in which quines or at least the though patterns involved in
coming to terms with them have remarkably wide relevance is in certain
data structures. A linked list might be presented in a diagram such as
\begin{verbatim}
     ____________      ____________
     |  A  :  --+----> |  B  :  --+----> ...
     ------------      ------------
\end{verbatim}
where rather small blocks of memory are chained together. Well if you have
a long list each block uses space, and sometimes you might like a list where
all the cells hold the same value, or at least where there is a short
repetition along it. Self-reference may support this nicely!
\begin{verbatim}
     ____________
     |  A  :  --+----
     ------------   |
      ^             |
      L_____________J
\end{verbatim}
The example of a re-entrant data structure here is tiny and not really
very exciting but the general idea has many many real-world applications,
but describing and reasoning about such can be a lot harder than
working with data that is not looped up. It may be easy to imagine that
if it was expressed using a programming system that supported this sort
of thing that setting up the description of code that works with copies
of itself as its input is really very much like the reference shown
here from within a data structure to its start-point.

This direction is the start of a path that leads toward models of
programming where the distinction between code and data is sometimes
blurred, and thus rather than programs being constrained to
having a linear forn with lines of code apearing sequentially in
a file, the real versions of them can have cyclic references as in the
simple list data here, and they can involve multiple references to
a single shared part thereby saving any need to having that written out
multiple times. A later chapter explores some of what emerges in
such a world.

     ____________      ____________
